% =====================================================================
% BACKGROUND SECTION
% Drop in via % =====================================================================
% BACKGROUND SECTION
% Drop in via % =====================================================================
% BACKGROUND SECTION
% Drop in via % =====================================================================
% BACKGROUND SECTION
% Drop in via \input{background_section} from main.tex.
% Assumes acmart sigconf and the bib entries provided in references.bib.
% =====================================================================

\section{Background}
\label{sec:background}

\subsection{DREM and the IOHMM framework}

The Dynamic Regulatory Events Miner (DREM) was introduced by Ernst et al.~\cite{ernst2007dremoriginal} as a method for reconstructing dynamic gene regulatory networks (GRNs) from time-series expression data. The model is an Input-Output Hidden Markov Model (IOHMM) over the temporal expression profile of each gene. Hidden states correspond to nodes in the dynamic network; observations are the per-time-point expression values of a gene; and the transition probabilities at each node are produced by a logistic regression classifier whose inputs are protein-DNA interaction values (the static input layer). When a node has more than one outgoing transition, that node is a \emph{split node}: a point along the trajectory (in real sampling time, or in pseudotime when DREM is run on a pseudotime-ordered series) at which a co-expressed group of genes diverges into multiple paths. The TF annotations at each split are based on the assignment of genes to their most likely path, which is determined by inference over the entire model. The logistic regression at each split then identifies which transcription factors (TFs) most distinguish genes assigned to one outgoing path from genes assigned to the other~\cite{ernst2007dremoriginal, schulz2012drem2}.

DREM 2.0~\cite{schulz2012drem2} extended the original framework with packaged protein-DNA interaction data for several species, support for continuous binding values, an option to scale TF effects by their own expression at each time point, dynamic (time-point-specific) input layers, and integrated motif discovery. The interactive GUI for inspection of the learned model was present in the original DREM and received smaller modifications in 2.0. The two pieces of input that the user supplies for a typical run are (i) a time-series matrix of dynamic measurements indexed by gene and time point, and (ii) a static or dynamic prior file in three columns: \texttt{TF}, \texttt{Gene}, \texttt{Input}, where the \texttt{Input} column is the binding value used by the logistic regression classifier (binary or continuous).

The first input encodes \emph{what is changing} along the trajectory. The second input encodes \emph{which regulators are candidates to explain the change}. DREM was originally designed assuming the first input is gene expression; whether it remains functional when methylation is supplied as the dynamic substrate instead is one of the questions this work evaluates. In principle the framework is agnostic about what biological measurement appears in the first input as long as it varies along time or pseudotime, and about what the prior connects as long as the rows of the dynamic matrix line up with the entities the prior maps to. Although DREM 2.0 supports dynamic (time-point-specific) priors, the work here uses a single static prior throughout. This separation between the dynamic substrate and the static prior, together with a correspondence between the row identifiers of the dynamic matrix and the entities indexed in the prior, is what the modality extension exploits.

\subsection{Methylation dynamics in the BICAN brain atlas}

The BICAN U01 project assembled a multi-modal single-cell atlas of the developing human brain spanning mid-gestation through 7 months postnatal, profiling DNA methylation (mCG) and three-dimensional chromatin contacts (3C) jointly via snm3C-seq~\cite{bican2026manuscript}. The developmental BICAN data were integrated with previously collected adult brain data, so that each lineage spans from a fetal progenitor state through to the corresponding adult cell type. Within that atlas, lineage trajectories were constructed for major neural cell types, including medial ganglionic eminence (MGE)-derived parvalbumin (PV) interneurons (MGE\_PVALB), two medium spiny neuron (MSN) subtypes (MSN\_DRD1\_BACH2 and MSN\_DRD1\_EPHA4), the oligodendrocyte precursor cell (OPC) to oligodendrocyte (ODC) differentiation arc (OPC\_ODC), and an immature caudal ganglionic eminence-derived LAMP5 inhibitory neuron lineage (eCGE\_LAMP5, where CGE denotes the caudal ganglionic eminence and the \texttt{e} prefix denotes the immature stage).

Differentially methylated regions (DMRs) were called by methylpy pairwise between adjacent developmental stages along each lineage (3T $\rightarrow$ 1m $\rightarrow$ 4--7m $\rightarrow$ adult). For each pairwise comparison, methylpy reports the differential regions from the perspective of each side, labeling both directions as ``hypo'' relative to the higher-methylation side. The merged lineage BED used here (despite its \texttt{\_hypo\_merged} filename) is the union of both directions of those pairwise comparisons, so it captures regulatory elements that lose mCG and regulatory elements that gain mCG along the trajectory, rather than demethylation events alone.

The dynamic measurement of interest in this setting is methylation $\beta$ at each DMR, indexed by developmental-stage cluster ordered along the lineage trajectory. This is the natural input to a DREM-style IOHMM if one wants to ask which TFs drive the methylation transitions at lineage bifurcations. The challenge is that DREM's standard prior is a TF-to-gene table built by intersecting ChIP-seq peaks with promoter windows, which is appropriate when the dynamic substrate is gene expression but mismatched when the dynamic substrate is methylation: the regions where methylation is changing are the DMRs themselves, not the $\pm 2$ kb windows around transcription start sites. Resolving this mismatch is the substantive methodological problem the rest of this report addresses.

\subsection{ChIP-Atlas and the TF binding evidence base}

The TF binding evidence used to build the prior comes from ChIP-Atlas~\cite{oki2018chipatlas}, a public reanalysis pipeline that uniformly processes ChIP-seq, ATAC-seq, and DNase-seq experiments deposited in NCBI SRA and similar archives. For each experiment, ChIP-Atlas publishes a peak BED file at multiple stringency thresholds and an experiment metadata table linking each experiment accession (SRX, ERX, or DRX, depending on the archive in which the experiment was originally deposited) to a target antigen and cell type. The peak BEDs are the substrate from which the TF-DMR prior is constructed: the framework intersects per-experiment peaks with lineage-specific DMR BEDs to produce a (TF, DMR) graph weighted by which transcription factors show binding evidence in the genomic regions where methylation is dynamic along each lineage. ChIP-Atlas's metadata also requires curation, since the ``antigen'' column includes entries that are not sequence-specific transcription factors (epitope tags, oxidative modifications, GFP fusions, histone modifications); these typically reflect a different experimental approach in which the intended target is annotated in a different metadata column, rather than annotation errors. The accession set used here filters these out; the antigens that remain are predominantly sequence-specific transcription factors, together with a minority of non-sequence-specific chromatin-associated proteins (for example cohesin and Mediator subunits and SWI/SNF remodeler subunits) that the curation procedure retained on the strength of their ChIP-seq evidence. Section~\ref{sec:framework:tfset} enumerates this composition.

\subsection{BrainSpan as a brain-relevance filter}

ChIP-Atlas aggregates experiments across all available cell types and tissues, so a TF that has many peaks in the genome may nonetheless be biologically irrelevant in the brain. To restrict the prior to TFs that are operative in human neural tissue, this work uses the BrainSpan developmental transcriptome~\cite{miller2014brainspan} as a second filter. BrainSpan provides bulk RNA-seq from 524 samples spanning multiple brain regions and developmental stages from 8 post-conceptional weeks through adulthood, processed to a per-gene RPKM matrix. A TF passes the filter if its mean RPKM across the BrainSpan samples exceeds a threshold. The threshold is chosen so that borderline regulators of demonstrated brain importance are retained: REST, for example, has a BrainSpan mean RPKM of $1.29$, and the threshold used here is set just below such biologically anchored points so that REST and similar regulators are not excluded. The choice is discussed in Section~\ref{sec:framework}.

\subsection{Related approaches}

Several approaches use temporal expression data to reconstruct regulatory networks. Friedman's dynamic Bayesian networks~\cite{friedman1998dbn} provide a formally clean framework for learning network structure from time-series expression but do not natively incorporate a static TF-DNA prior, and the structure-learning scales poorly to the hundreds of TFs that typical regulatory analyses require. Network Component Analysis~\cite{liao2003nca} does combine static interaction data, in the form of a TF-promoter connectivity matrix derived from ChIP-chip or literature, with expression measurements, but its matrix-factorization decomposition treats the time index as exchangeable rather than ordered. DREM's IOHMM occupies a useful middle ground: it respects time ordering, scales to hundreds of TFs through the logistic regression at split nodes, and produces an interpretable graphical model with TF annotations attached to specific bifurcation events. The work presented here treats DREM's IOHMM as a fixed inferential engine and focuses entirely on the prior side of the input pair, swapping the gene-substrate prior for a DMR-substrate prior matched to the methylation dynamic measurement.

% =====================================================================
% End of background section
% =====================================================================
 from main.tex.
% Assumes acmart sigconf and the bib entries provided in references.bib.
% =====================================================================

\section{Background}
\label{sec:background}

\subsection{DREM and the IOHMM framework}

The Dynamic Regulatory Events Miner (DREM) was introduced by Ernst et al.~\cite{ernst2007dremoriginal} as a method for reconstructing dynamic gene regulatory networks (GRNs) from time-series expression data. The model is an Input-Output Hidden Markov Model (IOHMM) over the temporal expression profile of each gene. Hidden states correspond to nodes in the dynamic network; observations are the per-time-point expression values of a gene; and the transition probabilities at each node are produced by a logistic regression classifier whose inputs are protein-DNA interaction values (the static input layer). When a node has more than one outgoing transition, that node is a \emph{split node}: a point along the trajectory (in real sampling time, or in pseudotime when DREM is run on a pseudotime-ordered series) at which a co-expressed group of genes diverges into multiple paths. The TF annotations at each split are based on the assignment of genes to their most likely path, which is determined by inference over the entire model. The logistic regression at each split then identifies which transcription factors (TFs) most distinguish genes assigned to one outgoing path from genes assigned to the other~\cite{ernst2007dremoriginal, schulz2012drem2}.

DREM 2.0~\cite{schulz2012drem2} extended the original framework with packaged protein-DNA interaction data for several species, support for continuous binding values, an option to scale TF effects by their own expression at each time point, dynamic (time-point-specific) input layers, and integrated motif discovery. The interactive GUI for inspection of the learned model was present in the original DREM and received smaller modifications in 2.0. The two pieces of input that the user supplies for a typical run are (i) a time-series matrix of dynamic measurements indexed by gene and time point, and (ii) a static or dynamic prior file in three columns: \texttt{TF}, \texttt{Gene}, \texttt{Input}, where the \texttt{Input} column is the binding value used by the logistic regression classifier (binary or continuous).

The first input encodes \emph{what is changing} along the trajectory. The second input encodes \emph{which regulators are candidates to explain the change}. DREM was originally designed assuming the first input is gene expression; whether it remains functional when methylation is supplied as the dynamic substrate instead is one of the questions this work evaluates. In principle the framework is agnostic about what biological measurement appears in the first input as long as it varies along time or pseudotime, and about what the prior connects as long as the rows of the dynamic matrix line up with the entities the prior maps to. Although DREM 2.0 supports dynamic (time-point-specific) priors, the work here uses a single static prior throughout. This separation between the dynamic substrate and the static prior, together with a correspondence between the row identifiers of the dynamic matrix and the entities indexed in the prior, is what the modality extension exploits.

\subsection{Methylation dynamics in the BICAN brain atlas}

The BICAN U01 project assembled a multi-modal single-cell atlas of the developing human brain spanning mid-gestation through 7 months postnatal, profiling DNA methylation (mCG) and three-dimensional chromatin contacts (3C) jointly via snm3C-seq~\cite{bican2026manuscript}. The developmental BICAN data were integrated with previously collected adult brain data, so that each lineage spans from a fetal progenitor state through to the corresponding adult cell type. Within that atlas, lineage trajectories were constructed for major neural cell types, including medial ganglionic eminence (MGE)-derived parvalbumin (PV) interneurons (MGE\_PVALB), two medium spiny neuron (MSN) subtypes (MSN\_DRD1\_BACH2 and MSN\_DRD1\_EPHA4), the oligodendrocyte precursor cell (OPC) to oligodendrocyte (ODC) differentiation arc (OPC\_ODC), and an immature caudal ganglionic eminence-derived LAMP5 inhibitory neuron lineage (eCGE\_LAMP5, where CGE denotes the caudal ganglionic eminence and the \texttt{e} prefix denotes the immature stage).

Differentially methylated regions (DMRs) were called by methylpy pairwise between adjacent developmental stages along each lineage (3T $\rightarrow$ 1m $\rightarrow$ 4--7m $\rightarrow$ adult). For each pairwise comparison, methylpy reports the differential regions from the perspective of each side, labeling both directions as ``hypo'' relative to the higher-methylation side. The merged lineage BED used here (despite its \texttt{\_hypo\_merged} filename) is the union of both directions of those pairwise comparisons, so it captures regulatory elements that lose mCG and regulatory elements that gain mCG along the trajectory, rather than demethylation events alone.

The dynamic measurement of interest in this setting is methylation $\beta$ at each DMR, indexed by developmental-stage cluster ordered along the lineage trajectory. This is the natural input to a DREM-style IOHMM if one wants to ask which TFs drive the methylation transitions at lineage bifurcations. The challenge is that DREM's standard prior is a TF-to-gene table built by intersecting ChIP-seq peaks with promoter windows, which is appropriate when the dynamic substrate is gene expression but mismatched when the dynamic substrate is methylation: the regions where methylation is changing are the DMRs themselves, not the $\pm 2$ kb windows around transcription start sites. Resolving this mismatch is the substantive methodological problem the rest of this report addresses.

\subsection{ChIP-Atlas and the TF binding evidence base}

The TF binding evidence used to build the prior comes from ChIP-Atlas~\cite{oki2018chipatlas}, a public reanalysis pipeline that uniformly processes ChIP-seq, ATAC-seq, and DNase-seq experiments deposited in NCBI SRA and similar archives. For each experiment, ChIP-Atlas publishes a peak BED file at multiple stringency thresholds and an experiment metadata table linking each experiment accession (SRX, ERX, or DRX, depending on the archive in which the experiment was originally deposited) to a target antigen and cell type. The peak BEDs are the substrate from which the TF-DMR prior is constructed: the framework intersects per-experiment peaks with lineage-specific DMR BEDs to produce a (TF, DMR) graph weighted by which transcription factors show binding evidence in the genomic regions where methylation is dynamic along each lineage. ChIP-Atlas's metadata also requires curation, since the ``antigen'' column includes entries that are not sequence-specific transcription factors (epitope tags, oxidative modifications, GFP fusions, histone modifications); these typically reflect a different experimental approach in which the intended target is annotated in a different metadata column, rather than annotation errors. The accession set used here filters these out; the antigens that remain are predominantly sequence-specific transcription factors, together with a minority of non-sequence-specific chromatin-associated proteins (for example cohesin and Mediator subunits and SWI/SNF remodeler subunits) that the curation procedure retained on the strength of their ChIP-seq evidence. Section~\ref{sec:framework:tfset} enumerates this composition.

\subsection{BrainSpan as a brain-relevance filter}

ChIP-Atlas aggregates experiments across all available cell types and tissues, so a TF that has many peaks in the genome may nonetheless be biologically irrelevant in the brain. To restrict the prior to TFs that are operative in human neural tissue, this work uses the BrainSpan developmental transcriptome~\cite{miller2014brainspan} as a second filter. BrainSpan provides bulk RNA-seq from 524 samples spanning multiple brain regions and developmental stages from 8 post-conceptional weeks through adulthood, processed to a per-gene RPKM matrix. A TF passes the filter if its mean RPKM across the BrainSpan samples exceeds a threshold. The threshold is chosen so that borderline regulators of demonstrated brain importance are retained: REST, for example, has a BrainSpan mean RPKM of $1.29$, and the threshold used here is set just below such biologically anchored points so that REST and similar regulators are not excluded. The choice is discussed in Section~\ref{sec:framework}.

\subsection{Related approaches}

Several approaches use temporal expression data to reconstruct regulatory networks. Friedman's dynamic Bayesian networks~\cite{friedman1998dbn} provide a formally clean framework for learning network structure from time-series expression but do not natively incorporate a static TF-DNA prior, and the structure-learning scales poorly to the hundreds of TFs that typical regulatory analyses require. Network Component Analysis~\cite{liao2003nca} does combine static interaction data, in the form of a TF-promoter connectivity matrix derived from ChIP-chip or literature, with expression measurements, but its matrix-factorization decomposition treats the time index as exchangeable rather than ordered. DREM's IOHMM occupies a useful middle ground: it respects time ordering, scales to hundreds of TFs through the logistic regression at split nodes, and produces an interpretable graphical model with TF annotations attached to specific bifurcation events. The work presented here treats DREM's IOHMM as a fixed inferential engine and focuses entirely on the prior side of the input pair, swapping the gene-substrate prior for a DMR-substrate prior matched to the methylation dynamic measurement.

% =====================================================================
% End of background section
% =====================================================================
 from main.tex.
% Assumes acmart sigconf and the bib entries provided in references.bib.
% =====================================================================

\section{Background}
\label{sec:background}

\subsection{DREM and the IOHMM framework}

The Dynamic Regulatory Events Miner (DREM) was introduced by Ernst et al.~\cite{ernst2007dremoriginal} as a method for reconstructing dynamic gene regulatory networks (GRNs) from time-series expression data. The model is an Input-Output Hidden Markov Model (IOHMM) over the temporal expression profile of each gene. Hidden states correspond to nodes in the dynamic network; observations are the per-time-point expression values of a gene; and the transition probabilities at each node are produced by a logistic regression classifier whose inputs are protein-DNA interaction values (the static input layer). When a node has more than one outgoing transition, that node is a \emph{split node}: a point along the trajectory (in real sampling time, or in pseudotime when DREM is run on a pseudotime-ordered series) at which a co-expressed group of genes diverges into multiple paths. The TF annotations at each split are based on the assignment of genes to their most likely path, which is determined by inference over the entire model. The logistic regression at each split then identifies which transcription factors (TFs) most distinguish genes assigned to one outgoing path from genes assigned to the other~\cite{ernst2007dremoriginal, schulz2012drem2}.

DREM 2.0~\cite{schulz2012drem2} extended the original framework with packaged protein-DNA interaction data for several species, support for continuous binding values, an option to scale TF effects by their own expression at each time point, dynamic (time-point-specific) input layers, and integrated motif discovery. The interactive GUI for inspection of the learned model was present in the original DREM and received smaller modifications in 2.0. The two pieces of input that the user supplies for a typical run are (i) a time-series matrix of dynamic measurements indexed by gene and time point, and (ii) a static or dynamic prior file in three columns: \texttt{TF}, \texttt{Gene}, \texttt{Input}, where the \texttt{Input} column is the binding value used by the logistic regression classifier (binary or continuous).

The first input encodes \emph{what is changing} along the trajectory. The second input encodes \emph{which regulators are candidates to explain the change}. DREM was originally designed assuming the first input is gene expression; whether it remains functional when methylation is supplied as the dynamic substrate instead is one of the questions this work evaluates. In principle the framework is agnostic about what biological measurement appears in the first input as long as it varies along time or pseudotime, and about what the prior connects as long as the rows of the dynamic matrix line up with the entities the prior maps to. Although DREM 2.0 supports dynamic (time-point-specific) priors, the work here uses a single static prior throughout. This separation between the dynamic substrate and the static prior, together with a correspondence between the row identifiers of the dynamic matrix and the entities indexed in the prior, is what the modality extension exploits.

\subsection{Methylation dynamics in the BICAN brain atlas}

The BICAN U01 project assembled a multi-modal single-cell atlas of the developing human brain spanning mid-gestation through 7 months postnatal, profiling DNA methylation (mCG) and three-dimensional chromatin contacts (3C) jointly via snm3C-seq~\cite{bican2026manuscript}. The developmental BICAN data were integrated with previously collected adult brain data, so that each lineage spans from a fetal progenitor state through to the corresponding adult cell type. Within that atlas, lineage trajectories were constructed for major neural cell types, including medial ganglionic eminence (MGE)-derived parvalbumin (PV) interneurons (MGE\_PVALB), two medium spiny neuron (MSN) subtypes (MSN\_DRD1\_BACH2 and MSN\_DRD1\_EPHA4), the oligodendrocyte precursor cell (OPC) to oligodendrocyte (ODC) differentiation arc (OPC\_ODC), and an immature caudal ganglionic eminence-derived LAMP5 inhibitory neuron lineage (eCGE\_LAMP5, where CGE denotes the caudal ganglionic eminence and the \texttt{e} prefix denotes the immature stage).

Differentially methylated regions (DMRs) were called by methylpy pairwise between adjacent developmental stages along each lineage (3T $\rightarrow$ 1m $\rightarrow$ 4--7m $\rightarrow$ adult). For each pairwise comparison, methylpy reports the differential regions from the perspective of each side, labeling both directions as ``hypo'' relative to the higher-methylation side. The merged lineage BED used here (despite its \texttt{\_hypo\_merged} filename) is the union of both directions of those pairwise comparisons, so it captures regulatory elements that lose mCG and regulatory elements that gain mCG along the trajectory, rather than demethylation events alone.

The dynamic measurement of interest in this setting is methylation $\beta$ at each DMR, indexed by developmental-stage cluster ordered along the lineage trajectory. This is the natural input to a DREM-style IOHMM if one wants to ask which TFs drive the methylation transitions at lineage bifurcations. The challenge is that DREM's standard prior is a TF-to-gene table built by intersecting ChIP-seq peaks with promoter windows, which is appropriate when the dynamic substrate is gene expression but mismatched when the dynamic substrate is methylation: the regions where methylation is changing are the DMRs themselves, not the $\pm 2$ kb windows around transcription start sites. Resolving this mismatch is the substantive methodological problem the rest of this report addresses.

\subsection{ChIP-Atlas and the TF binding evidence base}

The TF binding evidence used to build the prior comes from ChIP-Atlas~\cite{oki2018chipatlas}, a public reanalysis pipeline that uniformly processes ChIP-seq, ATAC-seq, and DNase-seq experiments deposited in NCBI SRA and similar archives. For each experiment, ChIP-Atlas publishes a peak BED file at multiple stringency thresholds and an experiment metadata table linking each experiment accession (SRX, ERX, or DRX, depending on the archive in which the experiment was originally deposited) to a target antigen and cell type. The peak BEDs are the substrate from which the TF-DMR prior is constructed: the framework intersects per-experiment peaks with lineage-specific DMR BEDs to produce a (TF, DMR) graph weighted by which transcription factors show binding evidence in the genomic regions where methylation is dynamic along each lineage. ChIP-Atlas's metadata also requires curation, since the ``antigen'' column includes entries that are not sequence-specific transcription factors (epitope tags, oxidative modifications, GFP fusions, histone modifications); these typically reflect a different experimental approach in which the intended target is annotated in a different metadata column, rather than annotation errors. The accession set used here filters these out; the antigens that remain are predominantly sequence-specific transcription factors, together with a minority of non-sequence-specific chromatin-associated proteins (for example cohesin and Mediator subunits and SWI/SNF remodeler subunits) that the curation procedure retained on the strength of their ChIP-seq evidence. Section~\ref{sec:framework:tfset} enumerates this composition.

\subsection{BrainSpan as a brain-relevance filter}

ChIP-Atlas aggregates experiments across all available cell types and tissues, so a TF that has many peaks in the genome may nonetheless be biologically irrelevant in the brain. To restrict the prior to TFs that are operative in human neural tissue, this work uses the BrainSpan developmental transcriptome~\cite{miller2014brainspan} as a second filter. BrainSpan provides bulk RNA-seq from 524 samples spanning multiple brain regions and developmental stages from 8 post-conceptional weeks through adulthood, processed to a per-gene RPKM matrix. A TF passes the filter if its mean RPKM across the BrainSpan samples exceeds a threshold. The threshold is chosen so that borderline regulators of demonstrated brain importance are retained: REST, for example, has a BrainSpan mean RPKM of $1.29$, and the threshold used here is set just below such biologically anchored points so that REST and similar regulators are not excluded. The choice is discussed in Section~\ref{sec:framework}.

\subsection{Related approaches}

Several approaches use temporal expression data to reconstruct regulatory networks. Friedman's dynamic Bayesian networks~\cite{friedman1998dbn} provide a formally clean framework for learning network structure from time-series expression but do not natively incorporate a static TF-DNA prior, and the structure-learning scales poorly to the hundreds of TFs that typical regulatory analyses require. Network Component Analysis~\cite{liao2003nca} does combine static interaction data, in the form of a TF-promoter connectivity matrix derived from ChIP-chip or literature, with expression measurements, but its matrix-factorization decomposition treats the time index as exchangeable rather than ordered. DREM's IOHMM occupies a useful middle ground: it respects time ordering, scales to hundreds of TFs through the logistic regression at split nodes, and produces an interpretable graphical model with TF annotations attached to specific bifurcation events. The work presented here treats DREM's IOHMM as a fixed inferential engine and focuses entirely on the prior side of the input pair, swapping the gene-substrate prior for a DMR-substrate prior matched to the methylation dynamic measurement.

% =====================================================================
% End of background section
% =====================================================================
 from main.tex.
% Assumes acmart sigconf and the bib entries provided in references.bib.
% =====================================================================

\section{Background}
\label{sec:background}

\subsection{DREM and the IOHMM framework}

The Dynamic Regulatory Events Miner (DREM) was introduced by Ernst et al.~\cite{ernst2007dremoriginal} as a method for reconstructing dynamic gene regulatory networks (GRNs) from time-series expression data. The model is an Input-Output Hidden Markov Model (IOHMM) over the temporal expression profile of each gene. Hidden states correspond to nodes in the dynamic network; observations are the per-time-point expression values of a gene; and the transition probabilities at each node are produced by a logistic regression classifier whose inputs are protein-DNA interaction values (the static input layer). When a node has more than one outgoing transition, that node is a \emph{split node}: a point along the trajectory (in real sampling time, or in pseudotime when DREM is run on a pseudotime-ordered series) at which a co-expressed group of genes diverges into multiple paths. The TF annotations at each split are based on the assignment of genes to their most likely path, which is determined by inference over the entire model. The logistic regression at each split then identifies which transcription factors (TFs) most distinguish genes assigned to one outgoing path from genes assigned to the other~\cite{ernst2007dremoriginal, schulz2012drem2}.

DREM 2.0~\cite{schulz2012drem2} extended the original framework with packaged protein-DNA interaction data for several species, support for continuous binding values, an option to scale TF effects by their own expression at each time point, dynamic (time-point-specific) input layers, and integrated motif discovery. The interactive GUI for inspection of the learned model was present in the original DREM and received smaller modifications in 2.0. The two pieces of input that the user supplies for a typical run are (i) a time-series matrix of dynamic measurements indexed by gene and time point, and (ii) a static or dynamic prior file in three columns: \texttt{TF}, \texttt{Gene}, \texttt{Input}, where the \texttt{Input} column is the binding value used by the logistic regression classifier (binary or continuous).

The first input encodes \emph{what is changing} along the trajectory. The second input encodes \emph{which regulators are candidates to explain the change}. DREM was originally designed assuming the first input is gene expression; whether it remains functional when methylation is supplied as the dynamic substrate instead is one of the questions this work evaluates. In principle the framework is agnostic about what biological measurement appears in the first input as long as it varies along time or pseudotime, and about what the prior connects as long as the rows of the dynamic matrix line up with the entities the prior maps to. Although DREM 2.0 supports dynamic (time-point-specific) priors, the work here uses a single static prior throughout. This separation between the dynamic substrate and the static prior, together with a correspondence between the row identifiers of the dynamic matrix and the entities indexed in the prior, is what the modality extension exploits.

\subsection{Methylation dynamics in the BICAN brain atlas}

The BICAN U01 project assembled a multi-modal single-cell atlas of the developing human brain spanning mid-gestation through 7 months postnatal, profiling DNA methylation (mCG) and three-dimensional chromatin contacts (3C) jointly via snm3C-seq~\cite{bican2026manuscript}. The developmental BICAN data were integrated with previously collected adult brain data, so that each lineage spans from a fetal progenitor state through to the corresponding adult cell type. Within that atlas, lineage trajectories were constructed for major neural cell types, including medial ganglionic eminence (MGE)-derived parvalbumin (PV) interneurons (MGE\_PVALB), two medium spiny neuron (MSN) subtypes (MSN\_DRD1\_BACH2 and MSN\_DRD1\_EPHA4), the oligodendrocyte precursor cell (OPC) to oligodendrocyte (ODC) differentiation arc (OPC\_ODC), and an immature caudal ganglionic eminence-derived LAMP5 inhibitory neuron lineage (eCGE\_LAMP5, where CGE denotes the caudal ganglionic eminence and the \texttt{e} prefix denotes the immature stage).

Differentially methylated regions (DMRs) were called by methylpy pairwise between adjacent developmental stages along each lineage (3T $\rightarrow$ 1m $\rightarrow$ 4--7m $\rightarrow$ adult). For each pairwise comparison, methylpy reports the differential regions from the perspective of each side, labeling both directions as ``hypo'' relative to the higher-methylation side. The merged lineage BED used here (despite its \texttt{\_hypo\_merged} filename) is the union of both directions of those pairwise comparisons, so it captures regulatory elements that lose mCG and regulatory elements that gain mCG along the trajectory, rather than demethylation events alone.

The dynamic measurement of interest in this setting is methylation $\beta$ at each DMR, indexed by developmental-stage cluster ordered along the lineage trajectory. This is the natural input to a DREM-style IOHMM if one wants to ask which TFs drive the methylation transitions at lineage bifurcations. The challenge is that DREM's standard prior is a TF-to-gene table built by intersecting ChIP-seq peaks with promoter windows, which is appropriate when the dynamic substrate is gene expression but mismatched when the dynamic substrate is methylation: the regions where methylation is changing are the DMRs themselves, not the $\pm 2$ kb windows around transcription start sites. Resolving this mismatch is the substantive methodological problem the rest of this report addresses.

\subsection{ChIP-Atlas and the TF binding evidence base}

The TF binding evidence used to build the prior comes from ChIP-Atlas~\cite{oki2018chipatlas}, a public reanalysis pipeline that uniformly processes ChIP-seq, ATAC-seq, and DNase-seq experiments deposited in NCBI SRA and similar archives. For each experiment, ChIP-Atlas publishes a peak BED file at multiple stringency thresholds and an experiment metadata table linking each experiment accession (SRX, ERX, or DRX, depending on the archive in which the experiment was originally deposited) to a target antigen and cell type. The peak BEDs are the substrate from which the TF-DMR prior is constructed: the framework intersects per-experiment peaks with lineage-specific DMR BEDs to produce a (TF, DMR) graph weighted by which transcription factors show binding evidence in the genomic regions where methylation is dynamic along each lineage. ChIP-Atlas's metadata also requires curation, since the ``antigen'' column includes entries that are not sequence-specific transcription factors (epitope tags, oxidative modifications, GFP fusions, histone modifications); these typically reflect a different experimental approach in which the intended target is annotated in a different metadata column, rather than annotation errors. The accession set used here filters these out; the antigens that remain are predominantly sequence-specific transcription factors, together with a minority of non-sequence-specific chromatin-associated proteins (for example cohesin and Mediator subunits and SWI/SNF remodeler subunits) that the curation procedure retained on the strength of their ChIP-seq evidence. Section~\ref{sec:framework:tfset} enumerates this composition.

\subsection{BrainSpan as a brain-relevance filter}

ChIP-Atlas aggregates experiments across all available cell types and tissues, so a TF that has many peaks in the genome may nonetheless be biologically irrelevant in the brain. To restrict the prior to TFs that are operative in human neural tissue, this work uses the BrainSpan developmental transcriptome~\cite{miller2014brainspan} as a second filter. BrainSpan provides bulk RNA-seq from 524 samples spanning multiple brain regions and developmental stages from 8 post-conceptional weeks through adulthood, processed to a per-gene RPKM matrix. A TF passes the filter if its mean RPKM across the BrainSpan samples exceeds a threshold. The threshold is chosen so that borderline regulators of demonstrated brain importance are retained: REST, for example, has a BrainSpan mean RPKM of $1.29$, and the threshold used here is set just below such biologically anchored points so that REST and similar regulators are not excluded. The choice is discussed in Section~\ref{sec:framework}.

\subsection{Related approaches}

Several approaches use temporal expression data to reconstruct regulatory networks. Friedman's dynamic Bayesian networks~\cite{friedman1998dbn} provide a formally clean framework for learning network structure from time-series expression but do not natively incorporate a static TF-DNA prior, and the structure-learning scales poorly to the hundreds of TFs that typical regulatory analyses require. Network Component Analysis~\cite{liao2003nca} does combine static interaction data, in the form of a TF-promoter connectivity matrix derived from ChIP-chip or literature, with expression measurements, but its matrix-factorization decomposition treats the time index as exchangeable rather than ordered. DREM's IOHMM occupies a useful middle ground: it respects time ordering, scales to hundreds of TFs through the logistic regression at split nodes, and produces an interpretable graphical model with TF annotations attached to specific bifurcation events. The work presented here treats DREM's IOHMM as a fixed inferential engine and focuses entirely on the prior side of the input pair, swapping the gene-substrate prior for a DMR-substrate prior matched to the methylation dynamic measurement.

% =====================================================================
% End of background section
% =====================================================================
